\chapter{Methodology}
\label{chap:methodology}

This chapter presents the methodology employed throughout this research. The experimental work was conducted in a systematic, phased approach: establishing retrieval baselines (Section~\ref{sec:meth_baseline}), performing error analysis to identify failure patterns (Section~\ref{sec:meth_error}), implementing and adapting the Query2Doc query enhancement technique (Section~\ref{sec:meth_query2doc}), and conducting a comprehensive model comparison across ten open-source Large Language Models (Section~\ref{sec:meth_model_comparison}). All experimental setup details, including dataset characteristics and hardware configurations, are described in Section~\ref{sec:meth_dataset}. The results corresponding to each methodological phase are presented in Chapter~\ref{chap:results}.

%======================================================================
\section{Dataset and Experimental Setup}
\label{sec:meth_dataset}
%======================================================================

\subsection{Dataset: MIRACL Arabic}
\label{sec:meth_miracl}

All experiments in this thesis were conducted on the Arabic subset of the Multilingual Information Retrieval Across a Continuum of Languages (MIRACL) benchmark \cite{zhang_2023_miracl}. MIRACL was selected as the primary dataset for three reasons: (1) it is retrieval-focused with human-annotated relevance judgements, (2) it contains natural query-document mismatch characteristic of real-world Arabic information needs, and (3) it is the standard benchmark for multilingual retrieval research, enabling comparison with published results.

The MIRACL Arabic corpus consists of 2,061,414 passages derived from 656,982 Arabic Wikipedia articles. Each passage is identified by a document identifier in the format \texttt{X\#Y}, where \texttt{X} denotes the article and \texttt{Y} the passage number within that article. The development set, which was used for all experiments reported in this thesis, contains 2,896 queries with 29,197 human-annotated relevance judgements. Each judgement assigns a relevance grade ranging from 0 (not relevant) to 3 (highly relevant), provided by native Arabic speakers. The corpus is written exclusively in Modern Standard Arabic (MSA); consequently, the evaluation of dialectal Arabic queries falls outside the scope of this work.

The dataset was accessed directly from HuggingFace\footnote{\url{https://huggingface.co/datasets/miracl/miracl}} without local download, as both the corpus and pre-built retrieval indices were available through the Pyserini toolkit \cite{lin_2021_pyserini}.

\subsection{Hardware and Software Environment}
\label{sec:meth_hardware}

All experiments were conducted on Google Colab, utilising two GPU configurations depending on availability and model requirements:

\begin{itemize}
    \item \textbf{NVIDIA T4 (15\,GB VRAM):} Used for baseline experiments and models with fewer than 4 billion parameters in FP16 precision.
    \item \textbf{NVIDIA A100 (40\,GB VRAM):} Used for models requiring 4-bit quantisation or exceeding 7 billion parameters, and for batch sizes larger than 8.
\end{itemize}

The software environment consisted of Python 3.10, PyTorch 2.0+, the Hugging Face Transformers library (v4.30+), Pyserini (v0.22.1) for index access, and FAISS for dense vector search. For models exceeding GPU memory in native precision, 4-bit NormalFloat (NF4) quantisation was applied using the \texttt{bitsandbytes} library \cite{dettmers_2023_qlora}, following the QLoRA quantisation scheme.

\subsection{Evaluation Metrics}
\label{sec:meth_metrics}

Three standard information retrieval metrics were employed for evaluation, as defined in Section~\ref{sec:math_models}:

\begin{itemize}
    \item \textbf{NDCG@10} (Equation~\ref{eq:ndcg}): Measures the quality of ranking within the top-10 retrieved passages. This was treated as the primary metric throughout all experiments.
    \item \textbf{Recall@10} (Equation~\ref{eq:recall}): Measures the proportion of all relevant passages that appear in the top-10 results.
    \item \textbf{MRR} (Equation~\ref{eq:mrr}): Measures how quickly the first relevant passage is retrieved.
\end{itemize}

Additionally, Recall@100 was computed to assess overall retrieval coverage. All metrics were calculated using the \texttt{pytrec\_eval} library over the full development set of 2,896 queries.

\subsection{Experimental Pipeline Overview}
\label{sec:meth_pipeline}

The experimental pipeline follows a two-phase architecture, as illustrated in Figure~\ref{fig:pipeline_overview}. In the first phase (query enhancement), each query from the MIRACL development set is processed by a query enhancement module, which may leave the query unchanged (baseline) or augment it using an LLM-generated pseudo-document. In the second phase (retrieval and evaluation), the enhanced query is passed to a retrieval system, and the top-100 results are evaluated against the MIRACL relevance judgements.

% TODO: Add actual figure
\begin{figure}[htbp]
    \centering
    \fbox{\parbox{0.85\textwidth}{\centering\vspace{2em}
    \textbf{[Figure placeholder: Experimental Pipeline Flowchart]}\\[0.5em]
    \small User Query $\rightarrow$ Query Enhancement Module $\rightarrow$ Enhanced Query $\rightarrow$ Retriever (Dense / BM25) $\rightarrow$ Top-100 Results $\rightarrow$ Evaluation Metrics
    \vspace{2em}}}
    \caption{Overview of the experimental pipeline. The query enhancement module is the variable component across experiments; the retrieval and evaluation components remain fixed.}
    \label{fig:pipeline_overview}
\end{figure}

A key design decision was to separate query generation from retrieval evaluation into two independent notebooks. The query enhancement phase (which requires GPU for LLM inference) saves enhanced queries to disk as serialised Python objects. The evaluation phase then loads these pre-generated queries and performs retrieval and metric computation. This two-notebook workflow enabled efficient reuse of enhanced queries across multiple retrieval configurations (Dense and BM25) without regenerating them.

%======================================================================
\section{Baseline Implementation}
\label{sec:meth_baseline}
%======================================================================

Two retrieval baselines were established to provide reference performance levels against which all query enhancement experiments would be compared: a dense retrieval baseline using mDPR and a sparse retrieval baseline using BM25S. Both baselines used an \texttt{IdentityEnhancer} that returned queries unchanged.

\subsection{Dense Retrieval Baseline (mDPR)}
\label{sec:meth_dense_baseline}

The dense retrieval baseline employed Multilingual Dense Passage Retrieval (mDPR), described in Section~\ref{sec:mdpr}. The mDPR encoder (\texttt{castorini/mdpr-tied-pft-msmarco}) was loaded on GPU, and queries were encoded in batches of 64. The pre-built FAISS index for MIRACL Arabic (5.47\,GB, 768-dimensional embeddings) was accessed via Pyserini. For each query, the top-100 passages were retrieved using cosine similarity.

The decision to use mDPR rather than stronger retrieval models such as BGE-M3 (described in Section~\ref{sec:dense_retrieval}) was deliberate. mDPR was pre-trained on MS~MARCO (English) and was \emph{not} fine-tuned on the MIRACL Arabic corpus. This intentionally ``weaker'' baseline was chosen to provide greater headroom for demonstrating the impact of query enhancement techniques. Stronger models that were trained directly on MIRACL would leave less room for measurable improvement.

\subsection{Sparse Retrieval Baseline (BM25S)}
\label{sec:meth_bm25_baseline}

The sparse retrieval baseline employed BM25S \cite{bm25s_2024}, a pure-Python implementation of the BM25 algorithm (described in Section~\ref{sec:sparse_retrieval}). BM25S was selected over the Pyserini-based BM25 implementation after encountering a blocking dependency conflict between Java 21 and Java 11 in the Google Colab environment. The BM25S library eliminates all Java dependencies while maintaining competitive performance.

The BM25S index was built from the MIRACL Arabic corpus with the following configuration:

\begin{itemize}
    \item \textbf{Tokenisation:} Arabic stemming using PyStemmer with NLTK Arabic stopwords (245+ words).
    \item \textbf{BM25 variant:} Lucene-style scoring with parameters $k_1 = 0.9$ and $b = 0.4$.
    \item \textbf{Index storage:} Google Drive ($\sim$0.5\,GB), loaded via symbolic links in Colab.
\end{itemize}

This configuration achieved 96\% of the official MIRACL Pyserini baseline performance, with the 4\% difference deemed acceptable given the benefits of a pure-Python pipeline that integrates directly with the query enhancement modules.

\subsection{Baseline Comparison Rationale}
\label{sec:meth_baseline_rationale}

Dense and sparse retrieval were tested \emph{separately} rather than combined into a hybrid system. This decision, confirmed during the project planning phase, was made to provide clearer insight into where query enhancement improvements originate. As discussed in Section~\ref{sec:ir_methods}, dense retrieval captures semantic similarity while sparse retrieval relies on exact term matching. Testing both baselines independently enables analysis of whether a given query enhancement technique benefits semantic matching, lexical matching, or both.

%======================================================================
\section{Error Analysis Methodology}
\label{sec:meth_error}
%======================================================================

Before implementing query enhancement techniques, a systematic error analysis was conducted on the dense retrieval baseline (Experiment~001) to identify the primary failure patterns and guide technique selection.

\subsection{Quantitative Analysis Framework}
\label{sec:meth_error_quant}

A quantitative analysis was performed over all 2,896 queries in the MIRACL Arabic development set. Queries were segmented into three performance categories based on their individual NDCG@10 scores:

\begin{itemize}
    \item \textbf{Failed:} NDCG@10 $< 0.3$
    \item \textbf{Mediocre:} $0.3 \leq$ NDCG@10 $< 0.7$
    \item \textbf{Successful:} NDCG@10 $\geq 0.7$
\end{itemize}

For each query, the following features were computed:

\begin{enumerate}
    \item \textbf{Query length} (in tokens), measured after whitespace tokenisation.
    \item \textbf{Score gap} between the top-ranked passage and the second-ranked passage, as a proxy for retriever confidence.
    \item \textbf{First relevant document rank}, indicating how deep the retriever must search before finding a relevant passage.
    \item \textbf{Coverage at multiple cut-offs} (top-10, top-20, top-50, top-100), measuring the proportion of queries with at least one relevant passage at each depth.
\end{enumerate}

The Pearson correlation coefficient was computed between query length and NDCG@10 to quantify the relationship between query verbosity and retrieval performance.

\subsection{Query Length Bucketing}
\label{sec:meth_error_buckets}

Queries were grouped into three length buckets to analyse the effect of query length on retrieval performance:

\begin{itemize}
    \item \textbf{Short:} 1--3 tokens (5.1\% of queries)
    \item \textbf{Medium:} 4--8 tokens (86.2\% of queries)
    \item \textbf{Long:} 9+ tokens (8.8\% of queries)
\end{itemize}

The mean NDCG@10 was computed for each bucket, and the performance gap between short and long queries was used to quantify the impact of information poverty on retrieval effectiveness.

\subsection{Failed Query Inspection}
\label{sec:meth_error_inspection}

The 20 worst-performing queries (NDCG@10 = 0.000, indicating zero relevant passages in the top-10) were manually inspected to identify recurring linguistic patterns. Each query was categorised by question type (factoid, definition, temporal), the presence of named entities, diacritics, technical terminology, and spelling variations. These observations informed the selection of the first query enhancement technique.

%======================================================================
\section{Query2Doc Implementation}
\label{sec:meth_query2doc}
%======================================================================

Based on the error analysis findings (presented in Section~\ref{sec:res_error}), Query2Doc \cite{wang_2023_query2doc} was selected as the first query enhancement technique. This section describes the Query2Doc approach, the modifications made for Arabic zero-shot application, and the implementation details.

\subsection{Query2Doc Technique}
\label{sec:meth_q2d_technique}

Query2Doc, introduced by Wang et al. \cite{wang_2023_query2doc}, uses an LLM to generate a pseudo-document---a short passage that attempts to answer the query. This pseudo-document is then concatenated with the original query to form an enhanced query. The rationale is that the pseudo-document introduces additional vocabulary, context, and semantic information that bridges the gap between the query and relevant passages.

As described in Section~\ref{sec:qe_techniques}, Query2Doc falls within the category of LLM-based generative query enhancement techniques. It differs from Hypothetical Document Embeddings (HyDE) \cite{gao_2022_precise} in a critical way: while HyDE uses the generated document \emph{alone} as the retrieval query, Query2Doc \emph{concatenates} the generated document with the original query, preserving the original query terms.

\subsection{Modifications from the Original Paper}
\label{sec:meth_q2d_modifications}

Several modifications were made to adapt Query2Doc for this research context:

\begin{enumerate}
    \item \textbf{Zero-shot prompting:} The original Query2Doc paper employed few-shot prompting with 4 demonstration examples using GPT-3 (\texttt{text-davinci-003}, 175B parameters). In this work, zero-shot prompting was adopted exclusively. This decision was driven by resource constraints (small open-source models have limited context windows for few-shot examples) and the desire to evaluate model capabilities without task-specific demonstrations.

    \item \textbf{Small open-source models:} Rather than using proprietary API-based models, all experiments utilised open-source models with 2--8 billion parameters that can run on Google Colab free-tier GPUs. This approach ensures reproducibility and eliminates API costs.

    \item \textbf{Arabic-only generation:} The system prompt explicitly instructs the model to respond in Arabic only, whereas the original paper targeted English queries.

    \item \textbf{Retriever-specific concatenation:} For dense retrieval, simple concatenation was used: $q_{\text{enhanced}} = q \oplus d_{\text{pseudo}}$, where $\oplus$ denotes string concatenation. The original paper recommends repeating the query $n=5$ times before concatenation for sparse retrieval to prevent term dilution, as discussed in Section~\ref{sec:meth_q2d_bm25}.
\end{enumerate}

\subsection{LLM Configuration}
\label{sec:meth_q2d_llm}

The initial Query2Doc experiment used Qwen~2.5~3B~Instruct (described in Section~\ref{sec:qwen25_3b}) as the generation model. The following system prompt was used across all experiments:

\begin{quote}
    \textit{``You are asked to write a passage that answers the given query. Do not ask the user for further clarification. Respond in Arabic only.''}
\end{quote}

The generation parameters were configured as follows:

\begin{table}[htbp]
    \centering
    \caption{LLM generation parameters for the initial Query2Doc experiment (Experiment~003).}
    \label{tab:q2d_params}
    \begin{tabular}{@{}ll@{}}
        \toprule
        \textbf{Parameter} & \textbf{Value} \\
        \midrule
        Maximum new tokens & 128 \\
        Temperature & 0.7 \\
        Top-$p$ (nucleus sampling) & 0.9 \\
        Batch size & 8 \\
        Precision & FP16 \\
        \bottomrule
    \end{tabular}
\end{table}

The maximum generation length of 128 tokens ($\sim$100 Arabic words) was selected as a balance between expansion quality and inference speed. Initial testing with 256 tokens showed diminishing returns in retrieval improvement while doubling generation time.

\subsection{Batch Processing Optimisation}
\label{sec:meth_q2d_batch}

A critical engineering challenge was generation throughput. Na\"ive sequential processing ($\sim$10 seconds per query) would require approximately 8 hours for the full development set. Three optimisations were implemented to achieve practical runtimes:

\begin{enumerate}
    \item \textbf{Batch processing:} Multiple queries were tokenised with left-padding and processed simultaneously through the LLM, yielding an 8$\times$ speedup.
    \item \textbf{Reduced token generation:} The maximum generation length was reduced from 256 to 128 tokens, yielding a 2$\times$ speedup with negligible quality loss.
    \item \textbf{Inference optimisations:} The model was set to evaluation mode (\texttt{model.eval()}), half-precision (\texttt{float16}) was used, and gradient computation was disabled (\texttt{torch.no\_grad()}).
\end{enumerate}

The combined effect was a 16$\times$ speedup, reducing the full-set generation time from $\sim$8 hours to $\sim$40 minutes on a T4 GPU. This enabled rapid iteration across multiple models and configurations.

\subsection{Dense Retrieval with Query2Doc}
\label{sec:meth_q2d_dense}

For dense retrieval experiments, the enhanced query was formed by simple concatenation:

\begin{equation}
    q_{\text{enhanced}} = q_{\text{original}} \oplus d_{\text{pseudo}}
    \label{eq:q2d_dense}
\end{equation}

The enhanced query was then encoded by the mDPR encoder and used for FAISS similarity search, following the same procedure as the baseline (Section~\ref{sec:meth_dense_baseline}). This approach leverages the semantic content of the pseudo-document to improve the query's embedding representation.

\subsection{BM25 Retrieval with Query2Doc}
\label{sec:meth_q2d_bm25}

For sparse retrieval, the same enhanced queries were evaluated against the BM25S index. Notably, the initial BM25 experiment (Experiment~004) used simple concatenation identical to the dense configuration, \emph{without} implementing the query repetition strategy recommended by Wang et al. \cite{wang_2023_query2doc}. The original paper recommends:

\begin{equation}
    q_{\text{enhanced}}^{\text{BM25}} = \underbrace{q \oplus q \oplus \cdots \oplus q}_{n \text{ times}} \oplus d_{\text{pseudo}}
    \label{eq:q2d_bm25}
\end{equation}

\noindent where $n=5$. This repetition preserves the importance of original query terms in the BM25 scoring function (Equation~\ref{eq:bm25}), preventing the pseudo-document from diluting the term weights of the original query. The omission of this strategy in the initial experiment, and its subsequent impact, is analysed in Section~\ref{sec:res_q2d_bm25}.

%======================================================================
\section{Model Comparison Methodology}
\label{sec:meth_model_comparison}
%======================================================================

Following the initial Query2Doc results with Qwen~2.5~3B, a comprehensive model comparison was conducted to identify the optimal LLM for Arabic query expansion. Ten open-source models were evaluated using the identical Query2Doc pipeline.

\subsection{Model Selection Criteria}
\label{sec:meth_model_selection}

Models were selected based on three criteria: (1) Arabic language support, either through dedicated Arabic training or strong multilingual capabilities; (2) parameter count within the 2--8 billion range to remain executable on Google Colab GPUs; and (3) open-source availability with permissive licences for research use. The selected models span three categories, all of which are described in detail in Section~\ref{sec:models_used}:

\begin{itemize}
    \item \textbf{Arabic-specialised:} Falcon-H1-3B (Section~\ref{sec:falcon_h1}), Jais-2-8B (Section~\ref{sec:jais2}), ALLaM-7B (Section~\ref{sec:allam}), and SILMA Kashif-2B (Section~\ref{sec:silma}).
    \item \textbf{Multilingual:} Qwen~2.5~3B (Section~\ref{sec:qwen25_3b}), Qwen~2.5~7B (Section~\ref{sec:qwen25_7b}), Qwen3-4B (Section~\ref{sec:qwen3_4b}), Qwen3-8B (Section~\ref{sec:qwen3_8b}), Gemma~3~4B (Section~\ref{sec:gemma3}), and Aya~Expanse~8B (Section~\ref{sec:aya}).
    \item \textbf{Experimental:} GPT-OSS-20B (Section~\ref{sec:gptoss}), a Mixture-of-Experts model included to test whether larger architectures provide benefits despite resource constraints.
\end{itemize}

Table~\ref{tab:model_configs} summarises the configuration used for each model.

\begin{table}[htbp]
    \centering
    \caption{Model configurations for the Query2Doc comparison experiments. All models used \texttt{max\_new\_tokens}~=~128 and \texttt{top\_p}~=~0.9 unless otherwise noted.}
    \label{tab:model_configs}
    \begin{tabular}{@{}llcccc@{}}
        \toprule
        \textbf{Model} & \textbf{Params} & \textbf{Precision} & \textbf{Temp.} & \textbf{Batch} & \textbf{GPU} \\
        \midrule
        Qwen 2.5 3B & 3B & FP16 & 0.7 & 8 & T4 \\
        SILMA 2B & 2B & FP16 & 0.1 & 16 & A100 \\
        Falcon-H1 3B & 3.15B & BF16 & 0.1 & 1\textsuperscript{*} & A100 \\
        Qwen3-4B & 4B & FP16 & 0.7 & 32 & A100 \\
        Gemma 3 4B & 4B & FP16 & 0.1 & 16 & A100 \\
        Qwen 2.5 7B & 7B & FP16 & 0.1 & 16 & A100 \\
        ALLaM 7B & 7B & NF4 & 0.7 & --- & A100 \\
        Jais-2 8B & 8B & BF16 & 0.1 & 8 & A100 \\
        Qwen3-8B & 8B & FP16 & 0.1 & 16 & A100 \\
        Aya Expanse 8B & 8B & NF4 & 0.1 & 8 & A100 \\
        GPT-OSS 20B & 20B & NF4 & 0.1 & --- & A100 \\
        \bottomrule
    \end{tabular}
    \\[0.5em]
    \raggedright\small\textsuperscript{*}Forced to batch size 1 due to architecture-specific batching bug (see Section~\ref{sec:meth_model_issues}).
\end{table}

\subsection{Standardised Evaluation Protocol}
\label{sec:meth_model_protocol}

To ensure fair comparison, all models were evaluated using an identical protocol:

\begin{enumerate}
    \item \textbf{Sanity check:} The first 5 queries were processed and the outputs were manually inspected to verify that the model (a) generates text in Arabic, (b) produces relevant content, and (c) does not exhibit degenerate behaviour (repetition, code output, or empty responses).
    \item \textbf{Full generation:} All 2,896 queries were processed using the Query2Doc pipeline, and the enhanced queries were saved to disk.
    \item \textbf{Dense evaluation:} Enhanced queries were evaluated against the mDPR dense retrieval baseline.
    \item \textbf{BM25 evaluation:} The same enhanced queries were evaluated against the BM25S sparse retrieval baseline.
\end{enumerate}

The work was divided between the two team members: one researcher tested Falcon-H1-3B, Jais-2-8B, ALLaM-7B, Qwen3-4B, and GPT-OSS-20B, while the other tested SILMA Kashif-2B, Qwen~2.5-7B, Qwen3-8B, Gemma~3~4B, and Aya~Expanse~8B. All experiments used the same codebase, system prompt, and evaluation pipeline to minimise confounding variables.

\subsection{Temperature Selection}
\label{sec:meth_temperature}

Temperature was not fixed at a single value across all models. Two considerations guided temperature selection:

\begin{enumerate}
    \item \textbf{Model-recommended values:} Some models specify recommended generation parameters. Falcon-H1 recommends temperature 0.1, noting that higher values ``largely drop performance.'' Qwen3 models require sampling (\texttt{do\_sample=True}) to avoid infinite repetitions with greedy decoding.
    \item \textbf{Empirical testing:} For SILMA Kashif-2B, two temperature values were compared (0.7 and 0.1). Temperature 0.1 yielded a 2.5\% improvement in NDCG@10 over 0.7, leading to its adoption as the default for subsequent experiments where no model-specific recommendation existed.
\end{enumerate}

\subsection{Model-Specific Technical Issues}
\label{sec:meth_model_issues}

Several models required specific engineering workarounds, documented here for reproducibility:

\begin{enumerate}
    \item \textbf{Falcon-H1-3B:} The hybrid Mamba2-Transformer architecture (described in Section~\ref{sec:falcon_h1}) contains a batching bug in the attention mask extension logic. When left-padded batches are used, the 4D causal attention mask is not extended as new tokens are generated, causing a shape mismatch. The workaround was to force \texttt{batch\_size=1}, resulting in significantly longer generation times but correct outputs.

    \item \textbf{Jais-2-8B:} This model requires BF16 precision rather than FP16 due to its Squared-ReLU activation function, which produces values that overflow FP16 range. Additionally, the \texttt{token\_type\_ids} parameter must be excluded from the model inputs, and the padding token must be explicitly set to the end-of-sequence token.

    \item \textbf{Qwen3-4B and Qwen3-8B:} The Qwen3 generation includes ``thinking'' tokens (internal chain-of-thought reasoning) that must be stripped from the output. The parameter \texttt{enable\_thinking=False} was passed during generation. Greedy decoding must not be used, as it causes infinite repetition loops.

    \item \textbf{ALLaM-7B:} As a preview-stage model, ALLaM exhibited a sentencepiece tokenizer artefact where the Unicode character ``\texttt{\char"2581}'' (lower one-eighth block) leaked into decoded outputs. This special character, used internally by the sentencepiece tokenizer to denote word boundaries, contaminated the pseudo-documents. The impact on retrieval performance is discussed in Section~\ref{sec:res_dropped}.

    \item \textbf{GPT-OSS-20B:} The Mixture-of-Experts architecture (32 experts, top-4 routing, described in Section~\ref{sec:gptoss}) proved impractical through 4-bit quantisation, exhibiting generation speeds approximately 70$\times$ slower than dense models of comparable active parameter count. Furthermore, its Harmony chat format required a forced-final-channel prefix to suppress English-language reasoning.
\end{enumerate}

\subsection{Quantisation Strategy}
\label{sec:meth_quantisation}

Models exceeding the available GPU memory in native precision were quantised using 4-bit NormalFloat (NF4) quantisation via the \texttt{bitsandbytes} library \cite{dettmers_2023_qlora}. The quantisation configuration was:

\begin{itemize}
    \item \textbf{Quantisation type:} NF4 (NormalFloat 4-bit)
    \item \textbf{Compute dtype:} BF16 (for matrix multiplication)
    \item \textbf{Double quantisation:} Enabled (quantises the quantisation constants for further memory savings)
\end{itemize}

This configuration was applied to Jais-2-8B (on T4 only; BF16 was used on A100), ALLaM-7B, Aya~Expanse~8B, and GPT-OSS-20B. Notably, Aya~Expanse~8B achieved the best overall performance despite being quantised to 4-bit, suggesting that model architecture and training data quality are more significant factors than numerical precision for this task.
